\section{Closed subiterations and countable determination}
\label{sec:iteration}

We now iterate the bridge. The main point is that a pure state on the
resulting algebra is uniquely determined on a suitable separable
subalgebra. To prove this, we first show that one may retain any collection
of bridge steps containing the dependencies of each of its members,
and then perform all omitted steps afterward.

\subsection{States and traces along a chain}

A chain $(B_\alpha)_{\alpha\le\lambda}$ of unital inclusions is
\emph{continuous} if
$B_\gamma=\overline{\bigcup_{\alpha<\gamma}B_\alpha}$ at every nonzero
limit ordinal $\gamma\le\lambda$.

\begin{lemma}\label{na:unique-chain}
Let $(B_\alpha)_{\alpha\le\lambda}$ be a continuous chain whose
successor inclusions have property \textup{(U)}. Then every inclusion
$B_\alpha\subset B_\beta$, for $\alpha\le\beta\le\lambda$, has
\textup{(U)}. The unique lifts of equivalent pure states are equivalent.
\end{lemma}
\begin{proof}
At a successor, unique extensions compose, and the intermediate lift
is pure by \cref{na:unique-lift}. At a limit, two extensions agree on
each earlier algebra, hence on their dense union. Existence follows
from state extension. Transfinite induction gives the first assertion.
For the second, a unitary implementing equivalence in the original
algebra also conjugates the unique lifts; use
\cref{na:unitary-states}.
\end{proof}

\begin{lemma}\label{na:inequivalent-limit}
Let $(B_\alpha)_{\alpha\le\lambda}$ be continuous, with $\lambda$ a
nonzero limit ordinal. Suppose pure states $h,k$ of $B_\lambda$ restrict
to inequivalent pure states of $B_\alpha$ for all $\alpha$ in some tail
below $\lambda$. Then $h\not\sim k$.
\end{lemma}
\begin{proof}
Otherwise realize $h,k$ by unit vectors $\xi,\eta$ in one irreducible
representation $\pi$ of $B_\lambda$. For each earlier $\alpha$ in the
given tail, the cyclic subspaces
\[
 M_\alpha=\overline{\pi(B_\alpha)\xi},\qquad
 N_\alpha=\overline{\pi(B_\alpha)\eta}
\]
are reducing and carry inequivalent irreducible representations of
$B_\alpha$. The orthogonal projection onto $N_\alpha$, restricted to
$M_\alpha$, is an intertwiner, so it is zero. Thus
$\langle\pi(a)\xi,\eta\rangle=0$ for every $a\in B_\alpha$.
Continuity gives the same equality for every $a\in B_\lambda$,
contradicting cyclicity of $\xi$.
\end{proof}

We will also use two elementary permanence facts. A continuous union
of simple unital algebras is simple: a proper quotient is injective,
and hence isometric, on every earlier algebra, so it is isometric on
their dense union and its completion. Compatible faithful tracial
states remain faithful at the limit. Indeed, the extended trace's
null ideal has zero intersection with every earlier algebra; the
quotient by this ideal is again isometric on a dense union, so the
ideal is zero. At successors, simplicity follows from
\cref{na:simple-extension} whenever \textup{(U)} and a faithful
conditional expectation are available.

\subsection{The initial algebra and its subiterations}

Let
\[
 A_0=\bigotimes_{m=1}^{\infty}M_2,\qquad
 p_n=e_{11}^{\otimes n}\otimes1\quad(n\ge1),\qquad p_0=1.
\]
Here $A_0$ is the norm closure of its unital matrix subalgebras
$M_{2^n}$, with the compatible normalized matrix traces.

\begin{lemma}\label{na:base-peak}
The algebra $A_0$ is separable, simple and infinite-dimensional, its
trace is faithful, and $(p_n)$ peaks to a pure state $f_0$.
Moreover, $\tau(p_n)=2^{-n}$. Every unital $C^*$-algebra containing $A_0$
has no finite-dimensional irreducible representations.
\end{lemma}
\begin{proof}
Simplicity and faithfulness follow from the permanence facts above
applied to the matrix stages. Separability and infinite dimension are
immediate from those stages. A state taking value one on $p_n$ must
restrict to the vector state of the first basis vector on $M_{2^n}$.
These vector states are compatible, so they determine exactly one
state on the dense union and its completion. This proves peaking and
purity. The trace formula is the normalized matrix trace formula.

A nonzero finite-dimensional irreducible representation of an
algebra containing $A_0$ restricts to a unital representation of
$A_0$. Its kernel is zero by simplicity, which would embed the
infinite-dimensional algebra $A_0$ in a finite-dimensional algebra.
\end{proof}

Consider a continuous chain $(X_\beta)_{\beta\le\delta}$ starting
from $X_0=A_0$. At an \emph{active} index $j<\delta$, perform the
bridge of \cref{na:bridge} using $(p_n)$ and a decreasing sequence
of projections $(q_n^j)$ that peaks in $X_j$ to a state inequivalent
to the $p_n$-state, with
\[
 q_0^j=1,\qquad \tau(q_n^j)=2^{-n}.
\]
Write $v_j$ for its stable unitary, so
$X_{j+1}=C^*(X_j,v_j)$. At every other index set $X_{j+1}=X_j$.
The traces are the compatible traces supplied by the bridges.
By \cref{na:unique-chain}, $(p_n)$ peaks at every stage. For each
active $j$, $(q_n^j)$ peaks at every stage from $j$ onward, and its
state is equivalent to the $p_n$-state from stage $j+1$ onward.

For $T\subset\delta$, define the unital subalgebra
\[
 A[T]=C^*\bigl(A_0,\{v_j:j\in T,\ j\text{ active}\}\bigr)
 \subset X_\delta.
\]
Thus $A[\beta]=X_\beta$ when the ordinal $\beta$ is regarded as the
set of smaller ordinals. Every element of $X_\delta$ lies in $A[T]$
for some countable $T$: approximate it by a sequence of polynomials,
each involving finitely many stable unitaries.

Choose countable sets $I_j\subset j$ such that
$q_n^j\in A[I_j]$ for every $n$ at each active $j$, and set
$I_j=\varnothing$ otherwise. Call $T$ \emph{dependency-closed} if
\begin{equation}\label{eq:dependency-closed}
 j\in T\quad\Longrightarrow\quad I_j\subset T.
\end{equation}
Every countable set is contained in a countable dependency-closed
set: repeatedly adjoin the dependencies of its members and take the
union over the natural numbers.

\begin{lemma}\label{na:inherited}
For every dependency-closed $T\subset\delta$, the algebras
$A[T\cap\beta]$, $\beta\le\delta$, form a continuous bridge
iteration in the inherited order, with identity steps at indices
outside $T$. In particular, $A[T]$ is simple and has a faithful
tracial state. If $T$ is countable, $A[T]$ is separable.
\end{lemma}
\begin{proof}
Let $k\in T$ be active. Both sequences for its bridge lie in
$A[T\cap k]$, by \eqref{eq:dependency-closed}. The original
expectations onto their generated abelian algebras restrict to this
smaller base, since their ranges already lie in it.
The restriction property \cref{na:hnn-restriction} identifies the
algebra obtained by adjoining $v_k$ with that reduced HNN extension.

The sequences peak to pure states of the smaller base by
\cref{na:peaking-restriction}, applied to its inclusion in $X_k$.
That same lemma gives each of these particular states its unique
extension to $X_k$. If the two smaller states were equivalent,
\cref{na:unitary-states} and uniqueness would make their extensions
equivalent, contrary to the choice of the original active step.
The smaller base contains $A_0$, so \cref{na:base-peak} excludes
finite-dimensional irreducibles. All hypotheses of
\cref{na:bridge} therefore hold in the inherited iteration.

Continuity follows because a polynomial uses only finitely many
indices. Simplicity follows by induction from
\cref{na:simple-extension} and the limit observation above.
Faithfulness follows either from the restricted trace on $X_\delta$
or from the same permanence argument. Finally, a countable set of
stable unitaries together with the separable algebra $A_0$ generates
a separable algebra.
\end{proof}

\subsection{Moving omitted steps to the end}

The preceding lemma proves unique extension along the inherited
iteration. We need the stronger assertion that it holds from the
entire inherited algebra into the original result.

\begin{theorem}\label{na:closed-extension}
For every dependency-closed $S\subset\delta$, the inclusion
$A[S]\subset X_\delta$ has property \textup{(U)}.
\end{theorem}
\begin{proof}
We will insert the omitted indices, in their original order, after
$A[S]$. Recall that the ordinal $j$ is the set of indices below $j$,
whereas $j+1$ also contains the index $j$ itself.
Fix an active $j\notin S$, and for
$j+1\le\beta\le\delta$ set
\begin{equation}\label{eq:reordered-bases}
 O_\beta=A[j\cup(S\cap\beta)],\qquad
 R_\beta=A[(j+1)\cup(S\cap\beta)].
\end{equation}
These index sets are dependency-closed. Initially,
$O_{j+1}=X_j$ and $R_{j+1}=X_{j+1}$; afterward both chains
perform precisely the inherited steps with indices in $S$.

\smallskip\noindent
\emph{The omitted pair stays inequivalent on $O_\beta$.}
Both $(p_n)$ and $(q_n^j)$ peak on every $O_\beta$ by
\cref{na:peaking-restriction}, since they peak in $X_\delta$.
Their states are initially inequivalent. Suppose $k>j$ is an active
member of $S$, and they are still inequivalent on $O_k$.
The $q_n^k$-sequence also lies in $O_k$ and peaks there.
Each of these three particular states has a unique extension to
$X_k$, by \cref{na:peaking-restriction}; this does not use the
theorem being proved.

In $X_k$, the $q_n^j$-state is equivalent to the $p_n$-state,
because step $j$ has already been performed. The $q_n^k$-state is
inequivalent to that state, because step $k$ is active. Equivalence
on $O_k$ would pass to these particular unique extensions.
Consequently, on $O_k$ the $q_n^k$-state is equivalent to neither
the $p_n$-state nor the $q_n^j$-state. The $k$-bridge joins only
the classes of its designated pair, by \cref{na:bridge}. It therefore
does not join the two states of the omitted $j$-bridge.
At limits their inequivalence persists by
\cref{na:inequivalent-limit}. Transfinite induction proves the claim.

\smallskip\noindent
\emph{The omitted adjunction commutes past the inherited steps.}
Write $H_j$ for the reduced HNN adjunction with the fixed data
$(p_n),(q_n^j)$. We claim, respecting the trace and named generators,
\begin{equation}\label{eq:omitted-adjunction}
 R_\beta\cong H_j(O_\beta)
 \qquad(j+1\le\beta\le\delta).
\end{equation}
This is true initially. At an active $k\in S$ above $j$, the data
for both adjunctions lie in $O_k$. The expectation from
$R_k=H_j(O_k)$ onto $O_k$, followed by either $k$-data expectation,
is a trace-preserving expectation onto that data algebra. By
\cref{na:trace-expectation}, it agrees with the expectation supplied
there by the corresponding peaking sequence. The same observation
applies with $j$ and $k$ interchanged. Thus the reordering
isomorphism \cref{na:reorder} applies and gives
\[
 H_k\bigl(H_j(O_k)\bigr)
 \cong H_j\bigl(H_k(O_k)\bigr),
\]
proving the next instance of \eqref{eq:omitted-adjunction}.
Other successor steps are identities. Figure~\ref{fig:reordering}
displays the two orders at an active index $k\in S$.

\begin{figure}[htbp]
\centering
\begin{tikzpicture}[>=Stealth,every node/.style={font=\small}]
 \node (o) at (0,0) {$O_k$};
 \node (j) at (6.4,0) {$H_j(O_k)=R_k$};
 \node (k) at (0,-2) {$H_k(O_k)=O_{k+1}$};
 \node[align=center] (jk) at (6.4,-2)
 {$R_{k+1}\cong H_k(H_j(O_k))$\\[3pt]
  $\phantom{R_{k+1}}\cong H_j(H_k(O_k))$};
 \draw[->] (o) -- node[above] {adjoin $v_j$} (j);
 \draw[->] (o) -- node[left] {adjoin $v_k$} (k);
 \draw[->] (j) -- node[right] {adjoin $v_k$} (jk);
 \draw[->] (k) -- node[below] {adjoin $v_j$} (jk);
\end{tikzpicture}
\caption{Reordering two adjunctions whose data all lie in $O_k$.
The arrows are inclusions. Both paths give the same faithfully
tracial algebra, under an isomorphism fixing $O_k,v_j,v_k$.}
\label{fig:reordering}
\end{figure}

At a limit, the isomorphisms already constructed agree on overlaps
because they fix the named generators. The expectations from the
$O_\beta$ onto the two fixed abelian algebras generated by $(p_n)$
and $(q_n^j)$ also agree on overlaps, by
\cref{na:trace-expectation}. Applying \cref{na:hnn-limit} to these
fixed data identifies the closure of the preceding $H_j(O_\beta)$
with $H_j$ of the limit base, proving
\eqref{eq:omitted-adjunction} there.

At $\beta=\delta$ we have
$O_\delta=A[S\cup j]$ and
$R_\delta=A[S\cup(j+1)]$.
The two sequences still peak to inequivalent states on $O_\delta$,
by the first part of the proof. Thus \cref{na:bridge} applies to
\eqref{eq:omitted-adjunction}, giving \textup{(U)} for
\[
 A[S\cup j]\subset A[S\cup(j+1)].
\]
If $j\in S$ or $j$ is inactive, this inclusion is the identity.
The continuous chain $\bigl(A[S\cup j]\bigr)_{j\le\delta}$ starts
at $A[S]$ and ends at $X_\delta$; continuity follows because a
polynomial involves only finitely many indices outside $S$.
Now \cref{na:unique-chain} proves the theorem.
\end{proof}

\subsection{A separable algebra determines each pure state}

The separable pure-restriction argument below is closely related to
\cite[Proposition~2.3 and concluding Remark~1]{AkemannWassermannWeaver2010}.
Here we arrange that the separable algebra is dependency-closed;
\cref{na:closed-extension} then adds uniqueness of the full extension.

\begin{lemma}\label{na:countable-determination}
For every pure state $h$ of $X_\delta$, there is a countable
dependency-closed $S\subset\delta$ such that $h|_{A[S]}$ is pure
and $h$ is its unique state extension to $X_\delta$.
\end{lemma}
\begin{proof}
Construct increasing countable dependency-closed sets $S_n$, starting
with $S_0=\varnothing$. Choose a countable dense subset $D_n$ of
$A[S_n]$. For each $a\in D_n$ and each integer $m\ge1$,
\cref{na:excision} supplies a positive contraction $x_{a,m}$ in
$X_\delta$ satisfying
\[
 h(x_{a,m})=1,\qquad
 \bigl\|x_{a,m}(a-h(a)1)x_{a,m}\bigr\|<1/m.
\]
Each witness has countable support. Adjoin all these supports to
$S_n$ and close under dependencies to obtain $S_{n+1}$.

Set $S=\bigcup_n S_n$ and $B=A[S]$. Suppose
$h|_B=t h_1+(1-t)h_2$ for states $h_1,h_2$ and $0<t<1$.
Each $h_i$ takes value one on every witness $x_{a,m}$.
A positive contraction of state value one fixes the state's GNS
vector, so $h_i(a)=h_i(x_{a,m}a x_{a,m})$ and
$h_i(x_{a,m}^2)=1$. Hence
\[
 |h_i(a)-h(a)|
 =\bigl|h_i\bigl(x_{a,m}(a-h(a)1)x_{a,m}\bigr)\bigr|<1/m.
\]
Letting $m$ tend to infinity shows that both states agree with $h$
on every $D_n$. Their union is dense in $B$, so $h|_B$ is pure.
Its unique extension is $h$ by \cref{na:closed-extension}.
\end{proof}

\begin{corollary}\label{na:state-count}
Put $\mathfrak c=2^{\aleph_0}$. If $|\delta|\le\mathfrak c$, then
$X_\delta$ has at most $\mathfrak c$ pure states.
\end{corollary}
\begin{proof}
There are at most
\[
 \max(\aleph_0,|\delta|)^{\aleph_0}
 \le\mathfrak c^{\aleph_0}=\mathfrak c
\]
countable subsets of $\delta$. For each one, the corresponding
algebra $A[S]$ is separable and has at most $\mathfrak c$ states:
values on a countable dense subset determine a state. By
\cref{na:countable-determination}, every pure state of $X_\delta$
is the unique extension of one such restriction. There are at most
$\mathfrak c\cdot\mathfrak c=\mathfrak c$ possibilities.
\end{proof}
